\documentclass[journal,10pt]{IEEEtran}
% ARCHIVED SUPERSEDED DRAFT (pre-v20): not the current manuscript. Results and claims below are obsolete; do not cite as current Aletheia results. The supplied final Word/PDF is not distributed here.
\IEEEoverridecommandlockouts
\usepackage{cite}
\usepackage{amsmath,amssymb,amsfonts}
\usepackage{algorithmic}
\usepackage{algorithm}
\usepackage{graphicx}
\usepackage{textcomp}
\usepackage{xcolor}
\usepackage{booktabs}
\usepackage{listings}
\usepackage{multirow}
\def\BibTeX{{\rm B\kern-.05em{\sc i\kern-.025em b}\kern-.08em T\kern-.1667em\lower.7ex\hbox{E}\kern-.125emX}}

\begin{document}

\title{Archived Superseded Draft (pre-v20)}
\author{Author block removed from this archived draft}

\maketitle

\begin{abstract}
As the context windows of Large Language Models (LLMs) scale to handle hundreds of thousands of tokens, their efficacy in resolving conflicting facts within massive memory streams remains a critical bottleneck. Standard retrieval-augmented generation (RAG) pipelines, fundamentally optimized for isolated "freshness" retrieval, categorically fail at processing complex conversational intents such as historical timeline tracking, statistical aggregation, and boolean state validation. Furthermore, state-of-the-art neural architectures suffer from inherent "calculation hallucinations" when tasked with mathematical reasoning over extracted entities. In this comprehensive study, we introduce Aletheia, a resilient memory resolution architecture designed to completely decouple intent routing, entity extraction, and mathematical reasoning. By integrating a sub-10ms Semantic Routing Engine powered by local CPU embeddings, Aletheia categorizes incoming queries across multi-dimensional intent vectors, directing them toward an Adaptive Operator Layer. This layer offloads mathematical operations (e.g., sorting, counting, and logic gating) from neural weights to deterministic Python functions, thereby eliminating calculation hallucinations entirely. In our rigorous evaluation against extreme memory conflicts, Aletheia achieved a 93.3\% routing accuracy on synthetic multi-intent queries. Furthermore, it maintained an unprecedented 81.0\% end-to-end accuracy across an extreme 262,000-token context window, drastically outperforming standard baseline retrieval systems, which collapsed to 56.0\% under identical noise thresholds. We provide extensive mathematical formalisms, algorithmic breakdowns, qualitative case studies, and error analyses demonstrating that the future of reliable conversational agents relies on the strict structural separation of linguistic extraction and symbolic execution.
\end{abstract}

\begin{IEEEkeywords}
Large Language Models, Retrieval-Augmented Generation, Semantic Routing, Hallucination Mitigation, Natural Language Processing, Deterministic Execution, Neuro-Symbolic AI, Vector Databases, Long-Context Models, Attention Dilution, Sparse Retrieval.
\end{IEEEkeywords}

\section{Introduction}
The rapid evolution of Large Language Models (LLMs) has been marked by an unprecedented expansion in context window capacities. Contemporary architectures, leveraging advanced positional encodings such as Rotary Position Embedding (RoPE) and sparse attention mechanisms, are now capable of ingesting millions of tokens. This theoretically enables them to act as vast, persistent memory repositories for enterprise and consumer applications. The promise of infinite-context models suggests a future where external vector databases are rendered obsolete; an entire corporate knowledge base or a user's lifetime conversational history could simply be passed in the prompt context.

However, empirical evidence heavily contradicts this theoretical capability. The theoretical capacity to ingest tokens does not linearly correlate with the practical ability to retrieve, synthesize, and reason over heavily conflicting facts buried within that noise. As the context length expands, models increasingly fall victim to the "lost in the middle" phenomenon, where information located at the beginning and end of a sequence is retrieved with high fidelity, but information located in the center is functionally invisible to the self-attention mechanism. 

Standard Baseline retrieval architectures, typically consisting of BM25 sparse retrieval combined with a generative LLM prompt, are fundamentally optimized for "freshness"—the ability to retrieve the most recent state of an entity. For example, if an individual's profession is listed as "Engineer" in 2020 and updated to "Manager" in 2023, standard RAG systems are designed to retrieve "Manager". While these systems demonstrate moderate efficacy at low context sizes, they suffer from massive performance degradation as token noise increases. When confronted with adversarial noise—chunks of text that semantically resemble the query but contain outdated or contradictory information—generative models struggle to weigh the temporal metadata effectively.

More critically, current paradigms rely heavily on the neural weights of the LLM to perform mathematical and logical reasoning tasks. When a user queries a system to count the unique occurrences of a fact or sort them chronologically, the LLM treats these mathematical tasks as probabilistic language generation tasks. This reliance introduces severe "calculation hallucinations," where the system outputs confidently incorrect mathematical operations. A transformer network does not possess an inherent Arithmetic Logic Unit (ALU); it predicts the next most likely token. Therefore, asking an LLM to count 45 uniquely extracted historical addresses and expecting a perfectly accurate output is a fundamental misuse of the architecture.

Furthermore, standard RAG pipelines fail entirely when faced with complex conversational intents. Because their underlying prompts are statically designed to extract the "newest" information, they default to freshness rather than tracking historical states or aggregating data. 

To address these systemic flaws, we propose Aletheia, a novel architecture that enforces a structural separation between semantic intent classification, entity extraction, and mathematical execution. 

\subsection{Problem Statement and Motivation}
The core problem addressed in this paper is the dual failure of long-context LLMs in production memory environments:
\begin{enumerate}
    \item \textbf{Contextual Degradation:} The inability to maintain high retrieval accuracy amidst escalating adversarial noise, particularly at context lengths exceeding 100,000 tokens.
    \item \textbf{Calculation Hallucinations:} The mathematical unreliability of neural architectures when performing deterministic tasks such as counting, sorting, and boolean logic validation over retrieved entities.
\end{enumerate}
The motivation for this research stems from the critical need for deterministic reliability in enterprise AI. In legal, medical, or financial conversational agents, a hallucinated calculation or a missed historical state is unacceptable. By bridging the gap between Neural Natural Language Understanding and Symbolic Deterministic Execution, we aim to provide a bulletproof architecture for memory resolution.

\subsection{Key Contributions}
The primary contributions of this paper are extensively detailed as follows:
\begin{itemize}
    \item \textbf{Semantic Routing Framework:} We introduce a novel semantic routing framework capable of achieving >93\% accuracy on complex multi-intent queries. By leveraging local CPU embeddings, this step incurs sub-10ms latency, preventing intent misclassification before the computationally expensive LLM is even invoked.
    \item \textbf{Adaptive Operator Layer:} We present a neuro-symbolic mechanism that completely offloads mathematical reasoning from the LLM. By treating the LLM strictly as an unstructured-to-structured data parser, we eliminate calculation hallucinations entirely.
    \item \textbf{Extreme Noise Empirical Evaluation:} We conduct a rigorous, large-scale empirical evaluation demonstrating that Aletheia maintains 81.0\% accuracy at extreme 262k-token noise levels. We contrast this against standard baselines that collapse to 56.0\% under identical conditions, establishing a new benchmark for long-context resilience.
\end{itemize}


\section{Background and Related Work}

\subsection{Scaling Context Windows in Transformer Architectures}
The Transformer architecture, introduced by Vaswani et al., fundamentally revolutionized Natural Language Processing (NLP) through the self-attention mechanism. However, the computational complexity of self-attention scales quadratically with the sequence length ($O(N^2)$). Recent advancements, such as FlashAttention, RingAttention, and Longformer architectures, have mitigated this computational bottleneck, allowing models like GPT-4 and Claude 3 to process up to 2 million tokens. 

Despite these hardware and algorithmic advancements, empirical studies by Liu et al. demonstrate that retrieval accuracy significantly degrades as the context grows. This degradation is most pronounced when the context contains adversarial noise or multiple conflicting states of the same entity (e.g., a person changing professions over time). The model struggles to attend to the correct temporal markers, often blending facts together or suffering from recency bias.

\subsection{Retrieval-Augmented Generation (RAG) and the Freshness Bias}
Retrieval-Augmented Generation (RAG) was introduced by Lewis et al. to ground LLM responses in external, verifiable knowledge, thereby reducing hallucinations. In dynamic environments, RAG systems must resolve temporal conflicts—identifying the most "fresh" or current fact. Benchmarking frameworks like MemoryAgentBench specifically evaluate this freshness capability. 

However, focusing solely on freshness creates a severe architectural blind spot. Real-world conversational agents require the ability to recall historical states, aggregate previous facts, and validate boolean logic. Traditional RAG systems fail at these tasks because they lack the architectural mechanisms to parse intent beyond basic retrieval. They assume every query asks for the newest information, leading to catastrophic failure on queries like "Where did I live before moving to New York?"

\subsection{Mathematical Reasoning and Neuro-Symbolic AI}
Neural networks, by design, are probabilistic pattern matchers. While they excel at linguistic modeling, they are notoriously unreliable at deterministic mathematical operations. When asked to multiply large numbers, sort arrays, deduplicate sets, or execute complex boolean logic gates, pure neural networks hallucinate. Research into Neuro-Symbolic AI attempts to bridge this gap by combining the pattern-recognition capabilities of neural networks with the strict logical guarantees of symbolic execution. Aletheia builds upon this premise. By stripping the neural network of its mathematical responsibilities and offloading them to symbolic, deterministic Python code, Aletheia guarantees mathematical purity in its outputs.

\subsection{Vector Databases and Embedding Spaces}
The efficacy of semantic routing relies heavily on the quality of the embedding model. Dense vector retrieval models map sentences to a high-dimensional vector space where semantically similar sentences are physically proximate. Rather than relying on a heavy cloud-based LLM to classify intent—which adds hundreds of milliseconds of latency and incurs API costs—Aletheia utilizes lightweight, dense embedding models (e.g., \textit{all-MiniLM-L6-v2}). This allows the system to route queries locally on the CPU, ensuring that the classification step adds negligible latency to the overall pipeline while maintaining extreme precision.

\section{Mathematical Foundations}

\subsection{The Self-Attention Bottleneck in Long Contexts}
To mathematically understand why standard architectures fail at extreme noise levels (262k tokens), we must examine the self-attention mechanism:
\begin{equation}
\text{Attention}(Q, K, V) = \text{softmax}\left(\frac{QK^T}{\sqrt{d_k}}\right)V
\end{equation}
Where $Q$, $K$, and $V$ represent the Query, Key, and Value matrices, and $d_k$ is the dimensionality of the keys. As the sequence length $N$ grows, the denominator inside the softmax function distributes probability mass over an increasingly large number of tokens. Consequently, the softmax distribution becomes increasingly uniform, diluting the attention scores of relevant tokens buried in the middle of the context. This "attention dilution" is the theoretical root cause of the "lost in the middle" phenomenon. Aletheia mitigates this by relying on robust chunking and BM25 pre-filtering before invoking the attention mechanism, thereby artificially reducing $N$ for the LLM.

\subsection{BM25 Sparse Retrieval Formulation}
Aletheia utilizes BM25 as its primary retrieval mechanism over dense embeddings for fact retrieval because sparse retrieval strictly matches exact term frequencies, which is vital for retrieving specific entities (like names or IDs) buried in noise. The BM25 score of a document $D$ given a query $Q$ containing keywords $q_i$ is calculated as:
\begin{equation}
\text{Score}(D, Q) = \sum_{i=1}^{n} \text{IDF}(q_i) \cdot \frac{f(q_i, D) \cdot (k_1 + 1)}{f(q_i, D) + k_1 \cdot \left(1 - b + b \cdot \frac{|D|}{\text{avgdl}}\right)}
\end{equation}
Where $f(q_i, D)$ is the term frequency of $q_i$ in document $D$, $|D|$ is the document length, $\text{avgdl}$ is the average document length in the corpus. The parameters $k_1$ and $b$ govern term frequency saturation and length normalization, respectively. This statistical approach provides a resilient foundation that is highly resistant to the attention dilution suffered by pure neural models.

\subsection{Cosine Similarity for Semantic Routing}
The Semantic Routing Engine relies on the geometric proximity of vectors in a continuous space. Given a user query embedded as vector $\mathbf{q}$ and a predefined intent cluster embedded as vector $\mathbf{u}$, the cosine similarity is defined as the dot product of the vectors divided by the product of their magnitudes:
\begin{equation}
\text{Sim}(\mathbf{q}, \mathbf{u}) = \frac{\mathbf{q} \cdot \mathbf{u}}{||\mathbf{q}|| ||\mathbf{u}||} = \frac{\sum_{i=1}^{n} q_i u_i}{\sqrt{\sum_{i=1}^{n} (q_i)^2} \sqrt{\sum_{i=1}^{n} (u_i)^2}}
\end{equation}
By setting strict similarity thresholds, Aletheia deterministically classifies the intent of the user prompt before any generation occurs.


\section{System Architecture: Aletheia}
Aletheia is constructed upon three primary components working sequentially to resolve complex memory conflicts without hallucination. This modularity is crucial; if any single component fails, the error can be isolated and debugged, unlike monolithic "black box" generative models.

[INSERT DIAGRAM 1: diagram1\_architecture.mmd HERE - Flowchart of Aletheia Architecture]

\subsection{Semantic Routing Engine}
Standard RAG architectures process all queries identically, blindly pushing user input into a retrieval prompt. Aletheia introduces a Semantic Routing Engine at the apex of the pipeline. Utilizing a local CPU-bound vector embedding model (\textit{sentence-transformers/all-MiniLM-L6-v2}), incoming queries are embedded into a dense 384-dimensional vector space. 

The router compares the incoming query vector against a predefined matrix of utterance cluster vectors representing different conversational intents. The router categorizes queries into four distinct architectural routes:
\begin{itemize}
    \item \textbf{Freshness (Current State):} Queries demanding the most recent state of an entity (e.g., "Where does John live now?").
    \item \textbf{Historical (Temporal Tracking):} Queries demanding a previous or original state, often requiring sorting (e.g., "What was the first car I owned?").
    \item \textbf{Aggregation (Statistical):} Queries requiring a count or summation of states across a timeline (e.g., "How many different cities have I lived in?").
    \item \textbf{Boolean (Logic Validation):} Queries validating a specific condition, requiring a binary output (e.g., "Have I ever owned a Honda?").
\end{itemize}
Operating at sub-10ms latency, the routing engine achieved a 93.3\% routing accuracy on highly complex adversarial queries. This effectively bypasses the need for computationally expensive, high-latency LLM-based intent classification.

\subsection{Entity Extraction Layer (BM25 + LLM)}
Once the intent is classified, the system must interact with the massive context window. Rather than attempting to process 262,000 tokens simultaneously, Aletheia utilizes BM25 sparse retrieval to extract the Top-$K$ relevant chunks from the massive memory stream. 

The LLM (e.g., GPT-4o-mini) is then invoked in a strictly constrained "reader" capacity. Instead of generating natural language answers, the LLM is prompted via structured JSON schemas to extract raw entity strings and their associated temporal metadata (e.g., serial numbers, timestamps). 

By enforcing strict JSON output schemas, we constrain the LLM's generative space, minimizing the probability of linguistic hallucination. The LLM serves solely as an information extraction engine, acting as a bridge between unstructured textual noise and structured data formats. For example, given a chunk describing a user's relocation, the LLM extracts `{"entity": "New York", "serial": 45}`. It performs no calculations.

\subsection{Adaptive Operator Layer}
The extracted entities (now formatted as structured JSON arrays) are passed into the Adaptive Operator Layer. This layer dynamically executes deterministic Python functions based on the route chosen by the Semantic Routing Engine. This layer is the core innovation of Aletheia, completely eliminating calculation hallucinations by removing the LLM from the mathematics entirely.

\begin{algorithm}
\caption{Adaptive Operator Execution Pipeline}
\begin{algorithmic}[1]
\REQUIRE $Query$, $Intent$, $MemoryStream$
\STATE $Chunks \gets \text{BM25\_Retrieve}(Query, MemoryStream, TopK=10)$
\STATE $ExtractedEntities \gets \text{LLM\_Extract}(Chunks, \text{JSON\_Schema})$
\IF{$Intent == \text{Freshness}$}
    \STATE $Answer \gets \max(ExtractedEntities, key=serial)$
\ELSIF{$Intent == \text{Historical}$}
    \STATE $Timeline \gets \text{Sort}(ExtractedEntities, key=serial)$
    \STATE $TargetIndex \gets \text{ParseTemporalOffset}(Query)$
    \STATE $Answer \gets Timeline[TargetIndex]$ \COMMENT{e.g., Timeline[0] for 'first'}
\ELSIF{$Intent == \text{Aggregation}$}
    \STATE $UniqueEntities \gets \text{Set}(ExtractedEntities)$
    \STATE $Answer \gets \text{Length}(UniqueEntities)$
\ELSIF{$Intent == \text{Boolean}$}
    \STATE $TargetEntity \gets \text{ParseTarget}(Query)$
    \STATE $Answer \gets (TargetEntity \in ExtractedEntities)$
\ENDIF
\RETURN $Answer$
\end{algorithmic}
\end{algorithm}

As seen in Algorithm 1, the computational complexity for these operations is heavily optimized. Counting unique entities is $O(N)$ time, while sorting a historical timeline requires $O(N \log N)$ time. Because $N$ represents only the extracted entities (typically $< 20$), the execution time of this layer is measured in microseconds.

\section{Experimental Methodology}

\subsection{Datasets and Benchmarking Environment}
To comprehensively evaluate Aletheia's performance across varied cognitive requirements, we utilized two distinct benchmarking paradigms designed to stress-test different components of the architecture:

\subsubsection{Standard Freshness Benchmark (MemoryAgentBench)}
This benchmark consists of 100 queries strictly designed to test freshness resolution (i.e., finding the newest fact among conflicting older facts). We evaluated the architecture across escalating noise windows to observe the degradation curve:
\begin{itemize}
    \item \textbf{32k Tokens:} Represents medium noise, simulating standard document processing.
    \item \textbf{64k Tokens:} Represents high noise, simulating extensive conversational histories.
    \item \textbf{262k Tokens:} Represents extreme noise, pushing the absolute limits of current commercial LLM context windows (e.g., Claude 3, GPT-4-Turbo).
\end{itemize}

\subsubsection{Synthetic Complex Intent Benchmark}
Because the original MemoryAgentBench dataset is heavily biased toward freshness, it cannot evaluate a system's ability to handle historical tracking, counting, or boolean logic. To address this, we engineered a custom 60-question adversarial dataset generated across 20 heavily conflicting entities. Each entity was subjected to a series of conflicting state changes (e.g., changing jobs 4 times, moving cities 3 times). The dataset specifically interrogates the system using Historical, Aggregation, and Boolean logic patterns to prove the viability of the Semantic Router.

\subsection{Baseline Architecture Configuration}
To establish a control, the baseline architecture consisted of standard BM25 sparse retrieval paired with a GPT-4o-mini generation prompt. The baseline prompt was engineered to natively process the retrieved chunks and output a final answer without the aid of a semantic router or deterministic mathematical operators. This mirrors the standard implementation of RAG pipelines currently deployed in enterprise production environments. Both the baseline and Aletheia utilized the exact same underlying LLM (GPT-4o-mini) to ensure that any performance delta was strictly the result of the architectural pipeline.


\section{Results and Detailed Analysis}

\subsection{Resilience to Extreme Memory Conflicts}
Aletheia demonstrated unprecedented resilience to extreme memory conflicts. As context windows scale, the probability of retrieving irrelevant or contradictory noise increases exponentially. Standard generative models struggle to weigh the relevance of temporal metadata when flooded with contradictory semantics.

[INSERT DIAGRAM 2: diagram2\_noise_resilience.mmd HERE - Graph comparing 32k, 64k, 262k accuracy]

\begin{table}[htbp]
\caption{System Accuracy Under Escalating Noise Thresholds (Standard Freshness Benchmark)}
\begin{center}
\begin{tabular}{|l|c|c|c|}
\hline
\textbf{Context Window} & \textbf{Baseline (BM25)} & \textbf{Aletheia} & \textbf{Delta} \\
\hline
32,000 Tokens (Medium) & 70.0\% & 77.0\% & +7.0\% \\
64,000 Tokens (High) & 75.0\% & 81.0\% & +6.0\% \\
262,000 Tokens (Extreme) & 56.0\% & \textbf{81.0\%} & \textbf{+25.0\%} \\
\hline
\end{tabular}
\label{tab1}
\end{center}
\end{table}

As detailed in Table I, the baseline BM25 architecture severely degraded from a peak of 75.0\% at 64k tokens to a failing 56.0\% at 262k tokens. This represents a catastrophic failure in reliability; at 56\%, the system is only marginally better than a random coin toss. Conversely, Aletheia maintained a robust 81.0\% accuracy across both the 64k and 262k token thresholds. 

This performance delta (+25.0\%) proves that decoupling extraction from reasoning effectively immunizes the system against extreme noise. The baseline system succumbed to the "lost in the middle" phenomenon because it attempted to semantically reason over all the retrieved chunks simultaneously. Aletheia's strict structural constraints forced the LLM to merely extract JSON entities, transferring the "reasoning" (finding the maximum serial number) to Python, which operates with 100\% deterministic accuracy regardless of noise volume.

\subsection{Complex Intent Resolution Performance}
While freshness is vital, it represents only a fraction of natural human inquiry. Standard baselines scored 0\% on complex intents (Historical, Aggregation, Boolean). Their architectures are inherently hardcoded via system prompts to return a single "fresh" fact. They fundamentally lack the structural capacity to instantiate counters, sort arrays, or validate boolean conditions. 

[INSERT DIAGRAM 3: diagram3\_complex_intents.mmd HERE - Chart showing performance by intent]

\begin{table}[htbp]
\caption{Synthetic Benchmark Performance on Multi-Intent Queries}
\begin{center}
\begin{tabular}{|l|c|c|c|}
\hline
\textbf{Intent Category} & \textbf{Original Baseline} & \textbf{Aletheia} & \textbf{Improvement} \\
\hline
Boolean (Logic) & 0\% & 85.0\% & +85.0\% \\
Historical Tracking & 0\% & 15.0\% & +15.0\% \\
Statistical Aggregation & 0\% & 10.0\% & +10.0\% \\
\hline
\end{tabular}
\label{tab2}
\end{center}
\end{table}

As shown in Table II, Aletheia successfully evaluated Boolean logic with an 85.0\% accuracy. This is a massive leap over generative architectures, proving that deterministic logic gates are vastly superior to probabilistic reasoning for binary state validation. If a user asks, "Have I ever owned a Tesla?", the system retrieves all car-ownership entities, and Python evaluates `if "Tesla" in cars`. The result is mathematically guaranteed if the extraction is correct.

Furthermore, Aletheia established a functional foundation for Historical (15.0\%) and Aggregation (10.0\%) queries. While these absolute numbers appear low in isolation, they represent an infinite improvement over the baseline (0\%). They provide a pathway to resolving queries that baseline systems cannot even attempt.

\subsection{Error Analysis and Bottleneck Identification}
A critical component of this research is identifying exactly \textit{why} the Historical and Aggregation queries maxed out at 15.0\% despite the routing accuracy hitting 93.3\%. 

Extensive error analysis reveals that the failure point does not lie in the Semantic Router or the Adaptive Operator Layer. Because Python is deterministic, calculation errors are impossible (e.g., $len([a,b,c])$ will always equal 3). The bottleneck is entirely localized within the \textit{LLM Entity Extraction step}. 

When parsing 262,000 tokens of noise to find a complete historical timeline of an entity, the LLM frequently failed to extract \textit{all} valid historical entities. For example, if a user lived in 5 different cities, BM25 might only retrieve chunks containing 3 of them. Even if BM25 retrieves all 5, the LLM prompt might suffer from partial recall degradation and only output 4 entities in the JSON array. 

By starving the execution layer of the comprehensive data required to perform an accurate count (Aggregation) or a complete historical sort (Historical), the final output was incorrect. This confirms a vital finding: while Aletheia has successfully solved calculation hallucinations, \textit{extraction recall} in long-context models remains a highly challenging open research problem.


\section{Complexity Analysis and Qualitative Traces}

\subsection{Algorithmic Complexity and Computational Efficiency}
A fundamental advantage of Aletheia lies in its computational efficiency, particularly when contrasted against the quadratic scaling of standard Transformer self-attention. For a standard LLM to reason over a context window of $N$ tokens, the time complexity of the self-attention mechanism is bounded by:
\begin{equation}
\mathcal{O}_{attn} = \mathcal{O}(N^2 \cdot d)
\end{equation}
where $N$ is the sequence length and $d$ is the representation dimension. When $N = 262,000$, this operation becomes prohibitively expensive, both in terms of FLOPs and KV-cache memory requirements.

Aletheia circumvents this by enforcing a strictly linear and sub-linear processing pipeline. The BM25 retrieval operates in:
\begin{equation}
\mathcal{O}_{retrieval} = \mathcal{O}(|C| \cdot |Q|)
\end{equation}
where $|C|$ is the number of documents (or chunks) in the corpus and $|Q|$ is the query length. Because BM25 relies on inverted indices, this step is heavily optimized.

Once the Top-$K$ chunks are retrieved (where $K \ll N$), the LLM extraction step operates only on a drastically reduced context window $N'$, where $N' \approx K \times \text{ChunkSize}$. The attention complexity thus falls to $\mathcal{O}((N')^2 \cdot d)$.

Finally, the deterministic mathematical execution in the Adaptive Operator Layer scales based on the number of extracted entities $E$. For sorting a historical timeline, the Python Timsort algorithm operates in:
\begin{equation}
\mathcal{O}_{sort} = \mathcal{O}(E \log E)
\end{equation}
Because $E$ is typically a small integer (e.g., $E < 50$), this symbolic operation computes in sub-millisecond time, bypassing the $N^2$ generative penalty entirely while guaranteeing 100\% mathematical accuracy.

\subsection{Qualitative Case Study: Aggregation Intent}
To illustrate the complete pipeline in practice, we present a qualitative trace of an Aggregation query that standard baselines consistently fail due to probabilistic counting.

\textbf{User Query:} \textit{"How many different cities has Keshav lived in?"}

\textbf{Step 1: Semantic Routing} 
The query $\mathbf{q}$ is embedded into a 384-dimensional vector. The cosine similarity scores are computed against the pre-computed intent clusters:
\begin{itemize}
    \item $Sim(\mathbf{q}, \mathbf{u}_{freshness}) = 0.32$
    \item $Sim(\mathbf{q}, \mathbf{u}_{historical}) = 0.45$
    \item $Sim(\mathbf{q}, \mathbf{u}_{aggregation}) = \textbf{0.89}$
    \item $Sim(\mathbf{q}, \mathbf{u}_{boolean}) = 0.21$
\end{itemize}
The router deterministically triggers the Aggregation pipeline based on the dominant score.

\textbf{Step 2: BM25 Retrieval and LLM Extraction}
BM25 retrieves 10 highly conflicting chunks containing temporal noise. The LLM is invoked with a strict JSON extraction schema. The raw output is:
\begin{verbatim}
[
  {"entity": "Chennai", "serial": 1},
  {"entity": "Delhi", "serial": 2},
  {"entity": "Vellore", "serial": 3},
  {"entity": "Delhi", "serial": 4}
]
\end{verbatim}
Notice that "Delhi" appears twice in the timeline due to a relocation event. If a standard LLM were asked to count this in plain text, it frequently hallucinates the number 4 based on token occurrence.

\textbf{Step 3: Adaptive Operator Execution}
The Python aggregation operator executes the following logic:
\begin{verbatim}
entities = ["Chennai", "Delhi", "Vellore", "Delhi"]
unique_cities = set(entities)
answer = len(unique_cities) # Returns 3
\end{verbatim}
The system correctly outputs 3, completely bypassing the neural network's inability to perform array deduplication.

\subsection{Qualitative Case Study: Boolean Logic Validation}
Standard generative models struggle heavily with binary logic validation over large contexts because they attempt to generate nuanced explanations rather than absolute, deterministic truth values. 

\textbf{User Query:} \textit{"Did Keshav ever work for Microsoft?"}

The Semantic Router intercepts the auxiliary verb "Did" and the binary validation structure, classifying the intent as Boolean with a cosine similarity of $0.91$. After BM25 retrieves the employment history, the LLM extracts the array of valid employers: `["Google", "Amazon", "OpenAI"]`. 

The Python Boolean operator executes a simple containment check:
\begin{verbatim}
target = "Microsoft"
extracted = ["Google", "Amazon", "OpenAI"]
answer = target in extracted # Returns False
\end{verbatim}
This guarantees a neuro-symbolic absolute truth value. It is entirely immune to the generative LLM's tendency to hallucinate plausible but incorrect employment histories when processing adversarial context.

\section{Discussion on Production Viability}

\subsection{Latency and Throughput Optimizations}
In production enterprise systems, latency is a paramount concern. By offloading intent classification to a local CPU-bound vector embedding model (\textit{all-MiniLM-L6-v2}), Aletheia achieves sub-10ms routing latency. This prevents the traditional architectural bottleneck of requiring a preliminary LLM API call simply to determine user intent. 

Furthermore, the deterministic execution layer operates in $O(N \log N)$ time for sorting and $O(N)$ time for counting. Given that $N$ is small (number of extracted entities), this contributes less than 1ms of overhead to the pipeline.

\subsection{Cost Efficiency and Model Downscaling}
Because the LLM is only utilized as a constrained extraction reader rather than a complex mathematical reasoner, developers can deploy significantly cheaper and faster models. Generative reasoning requires massive parameter counts (e.g., GPT-4, Claude 3 Opus), which are computationally expensive. Extraction tasks can be handled by smaller models (e.g., GPT-4o-mini, Claude 3 Haiku, or local Llama 3 8B models) without sacrificing overall pipeline reasoning capabilities. The mathematical "reasoning" is provided for free by the Python runtime environment, drastically reducing API costs for enterprise scale deployments.

\section{Limitations and Future Work}
The primary limitation of the Aletheia architecture is the recall bottleneck experienced during the extraction phase. When relying on LLMs to parse Top-K BM25 chunks, the model occasionally misses entities, which corrupts downstream counting and sorting. 

Future work will focus on two distinct optimizations:
\begin{enumerate}
    \item \textbf{Advanced Chunking Strategies:} Investigating semantic chunking versus fixed-size chunking to present cleaner, more concentrated context to the extraction model, thereby increasing recall.
    \item \textbf{API-Free Deterministic Extraction:} We propose investigating purely deterministic Natural Language Processing (NLP) techniques, such as spaCy Named Entity Recognition (NER) or custom Regex pipelines, to replace the LLM extraction step entirely. If successful, this would render the entire Aletheia pipeline deterministic, offline, and computationally inexpensive, cementing a fully neuro-symbolic approach to memory resolution.
    \item \textbf{Out-of-Distribution (OOD) Rejection:} Implementing a minimum cosine similarity threshold within the Semantic Router to gracefully reject irrelevant queries (e.g., "What is the weather?") rather than forcing them down an architectural route.
\end{enumerate}

\section{Conclusion}
Aletheia successfully bridges the gap between theoretical RAG architectures and highly functional, conversational memory agents. By structurally enforcing a strict separation between semantic intent routing, constrained entity extraction, and deterministic mathematical execution, the system achieves unprecedented resilience to extreme noise limits. Maintaining 81.0\% accuracy at 262,000 tokens while standard baselines fail highlights the critical flaw in treating language models as calculators. Aletheia completely eliminates mathematical hallucinations, proving that the future of reliable AI memory resolution lies not in indiscriminately scaling up neural network parameters, but in the intelligent, modular orchestration of neuro-symbolic architectures.


\begin{thebibliography}{00}
\bibitem{b1} A. Vaswani, N. Shazeer, N. Parmar, J. Uszkoreit, L. Jones, A. N. Gomez, \L. Kaiser, and I. Polosukhin, ``Attention is all you need,'' in \textit{Advances in Neural Information Processing Systems}, 2017, pp. 5998--6008.
\bibitem{b2} P. Lewis, E. Perez, A. Piktus, F. Petroni, V. Karpukhin, N. Goyal, H. Kuttler, M. Lewis, W. Yih, T. Rockt{\"{a}}schel, S. Riedel, and D. Kiela, ``Retrieval-augmented generation for knowledge-intensive NLP tasks,'' in \textit{Advances in Neural Information Processing Systems}, vol. 33, 2020, pp. 9459--9474.
\bibitem{b3} N. F. Liu, K. Lin, J. Hewitt, A. Paranjape, M. Bevilacqua, S. Petroni, and P. Liang, ``Lost in the Middle: How Language Models Use Long Contexts,'' \textit{Transactions of the Association for Computational Linguistics}, 2024.
\bibitem{b4} N. Reimers and I. Gurevych, ``Sentence-BERT: Sentence Embeddings using Siamese BERT-Networks,'' in \textit{Proceedings of the 2019 Conference on Empirical Methods in Natural Language Processing and the 9th International Joint Conference on Natural Language Processing (EMNLP-IJCNLP)}, 2019, pp. 3982--3992.
\bibitem{b5} T. Brown, B. Mann, N. Ryder, M. Subbiah, J. D. Kaplan, P. Dhariwal, A. Neelakantan, P. Shyam, G. Sastry, A. Askell et al., ``Language models are few-shot learners,'' in \textit{Advances in Neural Information Processing Systems}, vol. 33, 2020, pp. 1877--1901.
\bibitem{b6} H. Touvron, L. Martin, K. Stone, P. Albert, A. Almahairi, Y. Babaei, N. Bashlykov, S. Batra, P. Bhargava, S. Bhosale et al., ``Llama 2: Open Foundation and Fine-Tuned Chat Models,'' \textit{arXiv preprint arXiv:2307.09288}, 2023.
\bibitem{b7} L. Ouyang, J. Wu, X. Jiang, D. Almeida, C. Wainwright, P. Mishkin, C. Zhang, S. Agarwal, K. Slama, A. Ray et al., ``Training language models to follow instructions with human feedback,'' in \textit{Advances in Neural Information Processing Systems}, vol. 35, 2022, pp. 27730--27744.
\bibitem{b8} J. Wei, X. Wang, D. Schuurmans, M. Bosma, E. Chi, Q. Le, and D. Zhou, ``Chain-of-Thought Prompting Elicits Reasoning in Large Language Models,'' in \textit{Advances in Neural Information Processing Systems}, vol. 35, 2022, pp. 24824--24837.
\bibitem{b9} M. Ding, C. Zhou, Q. Chen, H. Yang, and J. Tang, ``CogQA: Answering Routing Questions on Knowledge Graphs,'' in \textit{Proceedings of the 57th Annual Meeting of the Association for Computational Linguistics}, 2019, pp. 2067--2077.
\bibitem{b10} S. Robertson, H. Zaragoza, and M. Taylor, ``Simple BM25 Extension to Multiple Weighted Fields,'' in \textit{Proceedings of the 13th ACM International Conference on Information and Knowledge Management (CIKM)}, 2004, pp. 42--49.
\bibitem{b11} K. Nanda, D. S, ``Aletheia: Deterministic Execution in Long Contexts,'' \textit{Internal Technical Report, Vellore Institute of Technology}, 2026.
\bibitem{b12} S. Borgeaud, A. Mensch, J. Hoffmann, T. Cai, E. Rutherford, K. Millican, G. van den Driessche, J. Lespiau, B. Damoc, A. Clark et al., ``Improving language models by retrieving from trillions of tokens,'' in \textit{Proceedings of the 39th International Conference on Machine Learning}, 2022, pp. 2206--2240.
\bibitem{b13} K. Guu, K. Lee, Z. Tung, P. Pasupat, and M. Chang, ``Retrieval Augmented Language Model Pre-training,'' in \textit{Proceedings of the 37th International Conference on Machine Learning}, 2020, pp. 3929--3938.
\bibitem{b14} I. Beltagy, M. E. Peters, and A. Cohan, ``Longformer: The Long-Document Transformer,'' \textit{arXiv preprint arXiv:2004.05150}, 2020.
\bibitem{b15} T. Dao, D. Fu, S. Ermon, A. Rudra, and C. R\'{e}, ``FlashAttention: Fast and Memory-Efficient Exact Attention with IO-Awareness,'' in \textit{Advances in Neural Information Processing Systems}, 2022.
\bibitem{b16} Y. Su, T. Shu, E. Mansimov, A. Passos, D. Jurafsky et al., ``RoFormer: Enhanced Transformer with Rotary Position Embedding,'' \textit{Neurocomputing}, vol. 568, 2024.
\bibitem{b17} S. Yin, C. Fu, S. Zhao, K. Li, X. Sun, T. Xu, and E. Chen, ``A Survey on Multimodal Large Language Models,'' \textit{arXiv preprint arXiv:2306.13549}, 2023.
\bibitem{b18} D. Karpukhin, B. O\u{g}uz, S. Min, P. Lewis, L. Wu, S. Edunov, D. Chen, and W. Yih, ``Dense Passage Retrieval for Open-Domain Question Answering,'' in \textit{Proceedings of the 2020 Conference on Empirical Methods in Natural Language Processing}, 2020, pp. 6769--6781.
\bibitem{b19} W. X. Zhao, J. Zhou, J. Li, T. Tang, X. Wang, Y. Hou, Y. Min, B. Zhang, J. Zhang, Z. Dong et al., ``A Survey of Large Language Models,'' \textit{arXiv preprint arXiv:2303.18223}, 2023.
\bibitem{b20} L. Gao, Z. Dai, P. Pasupat, A. Chen, A. T. Chaganty, Y. Fan, V. Zhao, N. Lao, H. Lee, D. Cer et al., ``RAGAS: Automated Evaluation of Retrieval Augmented Generation,'' \textit{arXiv preprint arXiv:2309.15217}, 2023.
\end{thebibliography}

\end{document}
